\scanpage{62}
\mspage{II. 10}
% lines 1-3
trouve une bonne réduction en les \deleted{\unsure{tous points de}} ``places'' de $K'$ au dessus de la place de $K$ définie par $R$). On a le

% lines 4-10
\uline{Lemme} 1.12. Soit $V$ un anneau de val. discrète, \inserted{de corps résiduel $k$,} \deleted{\unsure{dont le corps}} $\ell$ un nb premier $\ne\operatorname{car} k$, supposons que $k$ satisfasse l'hyp. de I 6.2. \inserted{les notations sont celles de I 1.} Alors pour tte représentation \inserted{$\ell$-adique} essentiellement abélienne de $\pi=\operatorname{Gal}(\overline K/K)$, la restriction à $I(\ell)$ est essentiellement triviale.

% lines 11-12
Se déduit de I 6.2 comme \deleted{\unsure{l}} on vient de déduire \unsure{1.4|1.3} de 1.1. via 1.2.

\margin{1) Pour le cas des pts de car. $\ell$, cf feuille \unsure{``le pb des puissances entières''|``le pb des poids entiers''} p. 13}

% lines 14-21
\uline{Remarque} 1.13. \deleted{\unsure{a}} Le seul exemple connu où on sache vérifier l'hypothèse \deleted{\unsure{de sem}} de l'action \deleted{abélienne} \inserted{\uline{abélienne et}\ \uline{semi-simple}} d'un groupe de Galois sur \inserted{un} $T_\ell(A_{\overline K})$ est celui où $A$ est une variété ``de type CM'' au sens de Shimura-Taniyama [\quad]. \inserted{Nous voyons donc que} \deleted{Dans} \unsure{ce cas} (le corps de définition de $A$ étant quelconque) $A$ devient radiciellement isogène, sur une extension \deleted{finie} séparable finie convenable, à \unsure{une|un} schéma abélien définissable sur un corps fini (cas car. $p>0$), resp. sur l'anneau des entiers d'un corps de nombres (cas de la car. $p=0$). Ce dernier cas \inserted{tout au moins} était bien connu [\quad].

% final line
Il semble plausible qu'en caractéristique
% archive sheet number: 61
